\chapter{总结与展望}

\section{总结}

  本文围绕人车模拟器脚本开发效率低、易出错、调试繁琐、环境部署困难等工程痛点，完成了面向人车模拟器的快速开发验证插件系统的研究、设计、实现与测试工作。论文以 VS Code 扩展框架为基础，构建了覆盖编码 — 校验 — 调试 — 部署全流程的一站式开发辅助工具，主要成果如下：
  第一，完成了插件系统的整体架构设计。采用了分层加模块化的设计思路，清晰划分了表现层、业务层和数据层的职责边界，同时确立了模块化、可扩展、易用性、实时性和跨平台这五大核心设计原则。这样设计出来的系统结构清晰，不仅方便后期维护，也为后续的功能迭代打下了基础。
  第二，实现了五大核心功能模块。分别为代码智能补全、实时错误检测、多场景联动调试、一体化环境打包、状态栏交互，通过抽象语法树解析、Levenshtein 编辑距离、状态机参数计数、调试适配器协议、绿色便携打包等关键技术，解决了专用 API 无提示、错误无法提前发现、调试配置重复、环境迁移复杂等实际问题。
  第三，完成了全流程系统测试。测试结果表明，插件功能完整、运行稳定，能够显著提升脚本编写速度、降低 API 调用错误率、缩短调试与部署耗时，完全满足教学、科研与团队协作等真实场景需求，具备较高的实用性与推广价值。
  第四，总结出了一套可复用的垂直领域 IDE 插件开发方法。本文用到的分层架构、接口设计和具体实现思路，对自动驾驶仿真、机器人控制、工业自动化这些同类垂直领域的 SDK 辅助工具开发，也有一定的参考价值。

\section{展望}

  尽管插件已实现核心功能并达到预期目标，但受限于开发周期与测试场景，系统仍存在可优化空间，未来可从以下方向继续完善：
  第一，增强静态检测能力。在现有参数数量校验基础上，增加参数类型、返回值类型、取值范围、空值判断等更严格的语法与语义检查，进一步提升错误覆盖率与精准度。
  第二，支持自定义与私有 API 扩展。允许用户导入自定义接口配置文件，实现团队私有函数提示、项目级 API 文档联动，提升插件灵活性与适用范围。
  第三，提升远程调试与网络兼容性。优化弱网、跨设备、容器化环境下的调试连接稳定性，支持日志实时上传、远程断点、多机协同调试等高级能力。
  第四，轻量化与增量打包优化。采用差量打包、压缩优化、组件裁剪等策略减小环境包体积，提升分发、传输与解压速度，适配批量部署场景。
  第五，引入 AI 辅助开发能力。结合代码大模型实现自动补全、函数生成、注释生成、代码重构等智能化能力，从 “辅助提示” 向 “智能开发” 升级。
  未来，插件可进一步与仿真平台联动，实现脚本在线校验、一键运行、数据可视化、任务自动化等扩展功能，为人车模拟器与智能驾驶仿真开发提供更高效、更智能、更统一的工具支撑。
